\documentclass[11pt,a4paper]{book}

\usepackage[T1]{fontenc}
\usepackage{lmodern}
\usepackage{amsmath,amssymb,amsthm,mathtools,bm}
\usepackage[margin=2.6cm]{geometry}
\usepackage{enumitem}
\usepackage{array}
\usepackage[colorlinks=true,linkcolor=black,citecolor=black,urlcolor=black]{hyperref}

\numberwithin{equation}{chapter}
\setlist{itemsep=2pt,topsep=4pt}

% --- macros -----------------------------------------------------------
\newcommand{\dd}{\mathrm{d}}
\newcommand{\ii}{\mathrm{i}}
\newcommand{\ee}{\mathrm{e}}
\newcommand{\vk}{\bm{k}}
\newcommand{\vr}{\bm{r}}
\newcommand{\vp}{\bm{p}}
\newcommand{\sgn}{\operatorname{sgn}}
\newcommand{\Imag}{\operatorname{Im}}
\newcommand{\Real}{\operatorname{Re}}
\newcommand{\dsdO}{\frac{\dd\sigma}{\dd\Omega}}
\newcommand{\afs}{\alpha_{\mathrm{f}}}
\newcommand{\kB}{\kappa_{\mathrm B}}

\title{Inelastic Scattering:\\[2pt]
       Classical, Quantum and Relativistic Theory}
\author{Draft chapters for a physics text\\
        (companion to \emph{Scattering by Long-Range Forces})}
\date{}

\begin{document}
\frontmatter
\maketitle
\tableofcontents

\chapter*{Scope and conventions}
\addcontentsline{toc}{chapter}{Scope and conventions}

A collision is \emph{elastic} when the colliding objects emerge with the same
internal states and the same identities, so that only their relative motion
changes. It is \emph{inelastic} when kinetic energy is converted into, or
released from, some other form: internal excitation of one or both partners,
ionization, rearrangement or capture, radiation, or the creation of new
particles. This volume is the companion of the volume on elastic scattering and
follows the same rules of selection. Two broad families of methods are omitted
on purpose: treatments restricted to interactions of finite range (partial
waves with short-range potentials, scattering lengths, Breit--Wigner
resonances of a short-range well), and treatments in which the interaction is
expanded as a perturbation (the Born and distorted-wave Born approximations,
Fermi's golden rule, tree-level Feynman diagrams, Bethe--Heitler and
Klein--Nishina cross sections). What remains is:
\begin{itemize}
  \item the dynamics-independent framework: kinematics, thresholds, channels,
        unitarity and its consequences;
  \item classical and semiclassical treatments in which the long-range
        Coulomb (or polarization) force is handled exactly: energy transfer,
        stopping, capture, radiation, transitions at avoided crossings;
  \item exact structural results for long-range forces in the multichannel
        problem: threshold laws, Rydberg series, the Wannier law;
  \item relativistic inelastic processes that are non-perturbative in the
        strength of the field: pair creation in strong fields, supercritical
        Coulomb centres.
\end{itemize}
Where a standard result is derived perturbatively (for example the Bethe
stopping formula) it is quoted, flagged as such, and used only for comparison
with the non-perturbative limit.

Part~I (Chapters~1--2) sets up kinematics and the classical theory. Part~II
(Chapters~3--5) treats quantum inelastic scattering: multichannel theory,
long-range thresholds, and the semiclassical impact-parameter method.
Part~III (Chapters~6--8) treats relativistic inelastic scattering and closes
with a comparison of the three settings.

\paragraph{Conventions.}
\begin{itemize}
  \item A \emph{channel} is a specification of the species and internal states
        of the fragments; $i$ labels the initial channel and $f$ a final one.
        $\Delta E=\varepsilon_f-\varepsilon_i$ is the internal excitation energy
        ($\Delta E>0$ for excitation, $\Delta E<0$ for de-excitation or
        ``superelastic'' collisions), and $Q=-\Delta E$ is the Q-value.
  \item $\mu$ is the reduced mass, $E$ the CM kinetic energy of the relative
        motion before the collision, $v$ the relative speed, $k_i=\mu v/\hbar$.
  \item For a Coulomb interaction $V=\alpha/r$ with $\alpha=Z_1Z_2e^2$
        (Gaussian units), the Sommerfeld parameter is $\eta=\alpha/\hbar v$.
        Bohr's parameter is $\kB=2z e^2/\hbar v=2\eta$ for a projectile of
        charge $ze$ and a unit-charge target particle.
  \item $\omega$ is the angular frequency of an internal transition
        ($\hbar\omega=\Delta E$), $m$ the electron mass, and $N$ the number
        density of targets.
  \item An asterisk marks CM quantities, the subscript L the laboratory frame.
        In relativistic sections $c=1$ is stated explicitly where used.
\end{itemize}

\mainmatter

% =====================================================================
\part{Foundations}
% =====================================================================

\chapter{Inelastic Collisions: Definitions and Kinematics}

\section{Channels and cross sections}

A collision $a+A\to b+B$ is described by its \emph{channel}: the set of fragment
species and internal states. The initial channel $i$ is fixed by the
preparation; each final channel $f$ has its own cross section,
\begin{equation}\label{eq:sigfi-def}
  \dd\dot N_f=J\,N\,\frac{\dd\sigma_{fi}}{\dd\Omega}\,\dd\Omega ,
\end{equation}
exactly as for elastic scattering. The \emph{total inelastic cross section} is
$\sigma_{\mathrm{inel}}=\sum_{f\neq i}\sigma_{fi}$ and the total cross section is
$\sigma_{\mathrm{tot}}=\sigma_{\mathrm{el}}+\sigma_{\mathrm{inel}}$. Because the
final channels have different momenta, the relation between cross section and
amplitude contains a velocity ratio. Let the stationary wave function behave as
$\psi\to\ee^{\ii k_iz}\phi_i+\sum_f f_{fi}\,\ee^{\ii k_fr}\phi_f/r$, where
$\phi_n$ are the internal states and $k_f$ is the final relative wave number.
The incident current is $\hbar k_i/\mu$ and the current into $\dd\Omega$ in
channel $f$ is $(\hbar k_f/\mu)|f_{fi}|^2\dd\Omega$, so
\begin{equation}\label{eq:sigfi}
  \frac{\dd\sigma_{fi}}{\dd\Omega}=\frac{k_f}{k_i}\,|f_{fi}(\theta,\varphi)|^2 .
\end{equation}
The kinematic factor $k_f/k_i$ vanishes at threshold and will be the origin of
all threshold laws.

\section{Non-relativistic kinematics with internal excitation}

Energy conservation in the CM frame reads $E=E'+\Delta E$, where $E'$ is the
final kinetic energy. The final wave number is
\begin{equation}\label{eq:kf}
  k_f^2=k_i^2-\frac{2\mu\,\Delta E}{\hbar^2}.
\end{equation}
The process is possible only if $E\ge\Delta E$ (a CM threshold); in the lab the
threshold kinetic energy of the projectile is
$T_{\mathrm{th}}=\Delta E\,(m_1+m_2)/m_2$, larger because the motion of the
centre of mass cannot be converted. Let $\rho=m_1/m_2$, and let the target be at
rest in the lab. The lab angle and cross section follow from the elastic
formulas with $\rho$ replaced by
\begin{equation}\label{eq:rhoprime}
  \rho'=\frac{\rho}{\sqrt{1-\Delta E/E}},
\end{equation}
the ratio of the CM velocity to the outgoing CM speed of the projectile:
\begin{align}
  \tan\theta_{\mathrm L}&=\frac{\sin\theta^*}{\rho'+\cos\theta^*},\label{eq:inel-labtan}\\
  \left.\dsdO\right|_{\mathrm L}&=
     \frac{(1+\rho'^2+2\rho'\cos\theta^*)^{3/2}}{|1+\rho'\cos\theta^*|}
     \left.\dsdO\right|^*. \label{eq:inel-labxs}
\end{align}
The lab energy of the scattered projectile is
\begin{equation}\label{eq:E1prime}
  E_1'=\frac{E_{\mathrm L}}{(m_1+m_2)^2}\Bigl[m_1^2+m_2^2\bigl(1-\tfrac{\Delta E}{E}\bigr)
    +2m_1m_2\sqrt{1-\tfrac{\Delta E}{E}}\,\cos\theta^*\Bigr],
\end{equation}
so a spectrum of scattered projectiles at fixed lab angle shows a peak for every
accessible $\Delta E$. This is the basis of energy-loss spectroscopy.

\section{Relativistic kinematics: invariant mass and thresholds}

In this section $c=1$. For $p_1+p_2\to p_3+p_4$ the invariants
$s=(p_1+p_2)^2$ and $t=(p_1-p_3)^2$ are as in the elastic case, but now
$m_3,m_4$ differ from $m_1,m_2$. In the CM frame
\begin{equation}\label{eq:rel-pstar}
  p^*=\frac{\lambda^{1/2}(s,m_1^2,m_2^2)}{2\sqrt s},\qquad
  p_f^*=\frac{\lambda^{1/2}(s,m_3^2,m_4^2)}{2\sqrt s},
\end{equation}
with the K\"all\'en function
$\lambda(x,y,z)=x^2+y^2+z^2-2xy-2xz-2yz$. The reaction can proceed only if
$\sqrt s\ge\sum_fm_f$. With the target at rest, $s=m_1^2+m_2^2+2m_2E_{\mathrm L}$
and the threshold lab kinetic energy is
\begin{equation}\label{eq:Tth}
  T_{\mathrm{th}}=\frac{\bigl(\sum_fm_f\bigr)^2-(m_1+m_2)^2}{2m_2}.
\end{equation}
Examples: $p+p\to p+p+\pi^0$ gives $T_{\mathrm{th}}=2m_\pi+m_\pi^2/2m_p\approx280$~MeV;
$p+p\to p+p+p+\bar p$ gives $T_{\mathrm{th}}=6m_p\approx5.6$~GeV, whereas a collider
of two equal beams needs only $E_{\mathrm{beam}}=2m_p$ each. The CM angular cross
section for a two-body final state is, with $\mathcal M$ the invariant amplitude
of Chapter~7,
\begin{equation}\label{eq:rel-inel-xs}
  \left.\dsdO\right|^*=\frac{1}{64\pi^2s}\,\frac{p_f^*}{p^*}\,|\mathcal M|^2 .
\end{equation}

\section{Minimal momentum transfer and the adiabatic cutoff}\label{sec:qmin}

In an elastic collision the momentum transfer $q=|\vk_i-\vk_f|$ ranges from $0$
(forward) to $2k$. In an inelastic collision $k_f\ne k_i$ and
\begin{equation}\label{eq:qmin}
  \hbar q_{\min}=\hbar(k_i-k_f)=\frac{p_i^2-p_f^2}{p_i+p_f}
  \simeq\frac{\Delta E}{v}\qquad(\Delta E\ll E),
\end{equation}
since $p_i^2-p_f^2=2\mu\Delta E$. The forward direction is no longer reachable.
Relativistically, with energy transfer $\omega=\Delta E$ and longitudinal
momentum transfer $\omega/v$, the invariant momentum transfer is bounded by
\begin{equation}\label{eq:tmin}
  |t|_{\min}=\frac{\omega^2}{v^2}-\omega^2=\frac{\omega^2}{\beta^2\gamma^2}.
\end{equation}
This result is the first of the structural facts of this volume. The conjugate
variable to $q$ is the impact parameter, so \eqref{eq:qmin} translates into a
largest effective impact parameter,
\begin{equation}\label{eq:bmax}
  b_{\max}=\frac{\hbar}{q_{\min}}=\frac{\hbar v}{\Delta E}\qquad
  \Bigl(\text{relativistically }b_{\max}=\frac{\gamma\hbar v}{\Delta E}\Bigr),
\end{equation}
Bohr's \emph{adiabatic distance}: a distant passage lasts a time $\sim b/v$, and
if that is long compared with $1/\omega$ the target follows the perturbation
adiabatically and is left unexcited. For an infinite-range force the elastic
cross section diverges because $q_{\min}=0$, whereas the inelastic cross section
is cut off at $b_{\max}$ and is finite (typically with a logarithm). Inelasticity
supplies its own infrared cutoff.

\section{Observables: stopping power, straggling, mean free path}

A projectile traversing a medium of number density $N$ loses energy in discrete
events. Define, for a differential cross section $\dd\sigma(W)$ in the energy
transferred $W$,
\begin{equation}\label{eq:S-def}
  S=-\frac{\dd E}{\dd x}=N\!\int W\,\dd\sigma(W),\quad
  \Omega^2=\frac{\dd\langle(\Delta E)^2\rangle}{\dd x}=N\!\int W^2\,\dd\sigma(W),\quad
  \lambda_{\mathrm{mfp}}=\frac1{N\sigma_{\mathrm{inel}}}.
\end{equation}
$S$ is the stopping power, $\Omega^2$ the energy straggling and
$\lambda_{\mathrm{mfp}}$ the mean free path between inelastic events. The
stopping power weights each event by $W$, so it is dominated by the large
logarithmic region of intermediate $b$, while $\sigma_{\mathrm{inel}}$ is
dominated by distant, weak events.

\section{Problems}
\begin{enumerate}
  \item Derive \eqref{eq:rhoprime}--\eqref{eq:E1prime} from the velocity
        addition between the CM and lab frames.
  \item Show that \eqref{eq:tmin} follows from the invariant
        $t=(p_1-p_3)^2$ evaluated at $\theta^*=0$ for $\omega\ll E$.
  \item Evaluate $T_{\mathrm{th}}$ from \eqref{eq:Tth} for
        $e^-+\mathrm{H}\to e^-+\mathrm{H}^*(n=2)$ and for
        $p+p\to p+p+p+\bar p$, and compare the lab and collider requirements.
\end{enumerate}

\section*{Notes and references}
The kinematics of inelastic collisions is treated in \cite{LL3} (Chapter XVIII), \cite{MM} and \cite{BJ}; relativistic kinematics and phase space in \cite{BK,Hagedorn,PDG}. The adiabatic cutoff and its role in energy loss are due to Bohr \cite{Bohr13,Bohr48}. Stopping power, straggling and the Bethe--Born theory are reviewed in \cite{Fano63,Inokuti71,Sigmund} and in the review on the passage of particles through matter in \cite{PDG}.

% =====================================================================
\chapter{Classical Inelastic Scattering}
% =====================================================================

\section{Energy transfer in a Coulomb collision}\label{sec:impulse}

A heavy projectile of charge $ze$ and speed $v$ passes a free electron, at rest,
at impact parameter $b$. If the deflection of the projectile is negligible the
trajectory is a straight line and the transverse impulse delivered to the
electron is
\begin{equation}\label{eq:impulse}
  \Delta p_\perp=\int\frac{ze^2\,b\,\dd t}{(b^2+v^2t^2)^{3/2}}=\frac{2ze^2}{bv}.
\end{equation}
The longitudinal impulse integrates to zero. The energy transferred is
\begin{equation}\label{eq:Wb}
  W(b)=\frac{\Delta p_\perp^2}{2m}=\frac{2z^2e^4}{mv^2b^2},
\end{equation}
which reaches the largest possible transfer $W_{\max}=2mv^2$ at
$b_{\min}=ze^2/mv^2$. Since $b$ and $W$ are related monotonically,
\begin{equation}\label{eq:dsdW}
  \frac{\dd\sigma}{\dd W}=2\pi b\left|\frac{\dd b}{\dd W}\right|=\frac{2\pi z^2e^4}{mv^2W^2},
\end{equation}
the Rutherford law expressed in energy transfer. The energy loss per unit path in
a medium with $N_e$ electrons per unit volume is then
\begin{equation}\label{eq:S-free}
  S=N_e\!\int W\,\dd\sigma=\frac{4\pi z^2e^4N_e}{mv^2}\,\ln\frac{b_{\max}}{b_{\min}} .
\end{equation}
The upper limit is the adiabatic distance \eqref{eq:bmax}; the lower is the
distance at which the impulse approximation breaks down. A free electron
supplies no $b_{\max}$ (it is not adiabatic), so for a free target the logarithm
diverges, just as in elastic Coulomb scattering. Binding of the electron supplies
the cutoff, as the next section shows.

\section{A bound target: the classical oscillator and Bohr's stopping formula}

Model a bound electron as a harmonic oscillator of frequency $\omega$ with
charge $-e$. The force from the passing charge is linear in the oscillator
coordinate, so the energy transferred is given \emph{exactly} by the Fourier
transform of the force at $\omega$: for an oscillator initially at rest,
$\Delta E=|\tilde F(\omega)|^2/2m$. For the straight-line trajectory with
$F_\perp=ze^2b/(b^2+v^2t^2)^{3/2}$ and
$F_\parallel=ze^2vt/(b^2+v^2t^2)^{3/2}$ the transforms involve the modified Bessel
functions $K_0$, $K_1$. With $\xi=\omega b/v$ and an isotropic oscillator,
\begin{equation}\label{eq:DEosc}
  \Delta E(b)=\frac{2z^2e^4}{mv^2b^2}\,\xi^2\bigl[K_1^2(\xi)+K_0^2(\xi)\bigr].
\end{equation}
For $\xi\ll1$, $\xi K_1\to1$ and $\xi K_0\to0$, and \eqref{eq:DEosc} reduces to
the free-electron result \eqref{eq:Wb}. For $\xi\gg1$ both Bessel functions decay
like $\ee^{-\xi}$, so $\Delta E\propto\ee^{-2\omega b/v}$: the adiabatic
suppression of Section~\ref{sec:qmin}, now derived. Integrating over impact
parameters,
\begin{equation}\label{eq:Bohr-int}
  S=N\!\int2\pi b\,\dd b\,\Delta E(b)=\frac{4\pi z^2e^4NZ}{mv^2}\,
    \int_{\xi_{\min}}^{\infty}\xi\bigl[K_1^2(\xi)+K_0^2(\xi)\bigr]\dd\xi ,
\end{equation}
for $Z$ electrons per target atom and $\xi_{\min}=\omega b_{\min}/v$. For small
$\xi_{\min}$ the integral equals $\ln(2\ee^{-\gamma_E}/\xi_{\min})+O(\xi_{\min})$,
where $\gamma_E$ is Euler's constant and $2\ee^{-\gamma_E}=1.123$. Hence
\begin{equation}\label{eq:Bohr}
  S=\frac{4\pi z^2e^4NZ}{mv^2}\,\ln\frac{1.123\,mv^3}{z\,e^2\,\omega}.
\end{equation}
This is Bohr's classical stopping formula. For a real atom $\omega$ is replaced
by a mean frequency $\bar\omega$ defined by the oscillator-strength-weighted
average of $\ln\omega_j$. Its range of validity is set by the requirement that
$b_{\min}=ze^2/mv^2$ exceed the quantum length $\hbar/mv$ below which the
electron cannot be localized, that is
\begin{equation}\label{eq:kappaB}
  \kB=\frac{2ze^2}{\hbar v}\gg1 .
\end{equation}
The Bohr parameter $\kB$ is twice the Sommerfeld parameter. Fast light
projectiles have $\kB\ll1$, where the classical formula fails and Bethe's
(perturbative, hence quoted rather than derived) formula
$S=(4\pi z^2e^4NZ/mv^2)\ln(2mv^2/I)$ with mean excitation energy $I$ applies.
Section~\ref{sec:bloch} shows that a single function interpolates between the
two.

\section{Thomson's classical ionization cross section}

Treat the bound electron as free but requiring an energy transfer of at least
the ionization potential $I$. For an incident electron of kinetic energy $E$,
\eqref{eq:dsdW} reads $\dd\sigma/\dd W=\pi e^4/(EW^2)$. If the projectile and
target electron are treated as distinguishable and $W$ runs from $I$ to $E$,
\begin{equation}\label{eq:Thomson}
  \sigma_{\mathrm{ion}}=\frac{\pi e^4}{E}\Bigl(\frac1I-\frac1E\Bigr).
\end{equation}
Including exchange classically, that is summing the Rutherford terms for the two
outgoing electrons with $W\le E/2$ the energy of the slower one, the integrand is
$(\pi e^4/E)[1/W^2+1/(E-W)^2]$ and
\begin{equation}\label{eq:Thomson-ex}
  \sigma_{\mathrm{ion}}=\frac{\pi e^4}{E}\Bigl(\frac1I-\frac1{E-I}\Bigr).
\end{equation}
Both vanish \emph{linearly} at threshold: with $E=I+\varepsilon$,
$\sigma_{\mathrm{ion}}\simeq\pi e^4\varepsilon/I^3$. Section~\ref{sec:wannier}
shows that the true threshold law, which accounts for the long-range Coulomb
forces among the three bodies, has a different exponent.

\section{Capture by a long-range attraction}\label{sec:capture}

A particle in an attractive potential $V=-C/r^n$ with $n>2$ moves in the
effective potential $V_{\mathrm{eff}}=L^2/(2\mu r^2)-C/r^n$, which has a barrier
maximum at $r_m^{n-2}=nC\mu/L^2$ of height $V_m=(n-2)L^2/(2n\mu r_m^2)$. If the
energy exceeds the barrier the particle falls to the centre. With
$L^2=2\mu Eb^2$ the condition $V_m<E$ determines a critical impact parameter and
a capture cross section
\begin{equation}\label{eq:sigcap}
  \sigma_{\mathrm{cap}}(E)=\pi b_c^2=\pi\Bigl(\frac{n}{n-2}\Bigr)^{(n-2)/n}
     \Bigl(\frac{nC}{2E}\Bigr)^{2/n}\ \propto\ E^{-2/n}.
\end{equation}
Capture removes the particle from the elastic channel into a bound or reactive
channel, so it is an inelastic cross section, and it is finite even though the
force has infinite range. Two cases are important.
\begin{itemize}
  \item \textbf{Langevin capture} ($n=4$). A charge $ze$ induces a dipole
        $\alpha_dze/r^2$ in a neutral particle of polarizability $\alpha_d$,
        giving $V=-C/r^4$ with $C=\alpha_dz^2e^2/2$. Then
        $\sigma_L=2\pi\sqrt{C/E}=\pi ze\sqrt{2\alpha_d/E}$. The rate coefficient
        $k_L=v\sigma_L=2\pi ze\sqrt{\alpha_d/\mu}$ is independent of energy, which is why
        ion--molecule reaction rates are nearly temperature independent.
  \item \textbf{The limit $n\to2$.} The prefactor tends to 1 and
        $\sigma_{\mathrm{cap}}\to\pi C/E$, the capture cross section of the
        attractive inverse-square potential.
\end{itemize}
For a pure Coulomb attraction ($n=1$) the effective potential has no barrier
and classical two-body capture does not occur: capture into Coulomb bound states
requires an additional energy-loss mechanism (radiation or a third body).

\section{Classical bremsstrahlung and the soft-photon theorem}\label{sec:soft}

A charge that is accelerated radiates, so a Coulomb deflection is an inelastic
process. For non-relativistic motion the Larmor formula gives the radiated
energy spectrum of a single collision. In the low-frequency limit, where the
collision time is short compared with $1/\omega$, the spectrum depends only on the
net change of velocity $\Delta\bm v$ and not on the details of the orbit:
\begin{equation}\label{eq:soft-nr}
  \frac{\dd E_{\mathrm{rad}}}{\dd\omega}\bigg|_{\omega\to0}=
  \frac{2e^2}{3\pi c^3}\,|\Delta\bm v|^2 .
\end{equation}
This is exact for any trajectory. For an elastic Coulomb deflection through angle
$\Theta$, $|\Delta\bm v|^2=2v^2(1-\cos\Theta)$, so averaging over collisions
weights the radiation by the momentum-transfer cross section
$\sigma_{\mathrm{tr}}=\int(1-\cos\theta)\,\dd\sigma$:
\begin{equation}\label{eq:brems-sigtr}
  \frac{\dd\sigma_E}{\dd\omega}=\frac{4e^2v^2}{3\pi c^3}\,\sigma_{\mathrm{tr}} .
\end{equation}
For an electron on a nucleus of charge $Ze$, $\sigma_{\mathrm{tr}}=4\pi Z^2e^4
m^{-2}v^{-4}\ln(b_{\max}/b_{\min})$ with $b_{\min}=Ze^2/mv^2$, and the collision
duration condition gives $b_{\max}\approx v/\omega$:
\begin{equation}\label{eq:brems}
  \frac{\dd\sigma_E}{\dd\omega}=\frac{16}{3}\,\frac{Z^2e^6}{m^2c^3v^2}\,
  \ln\frac{b_{\max}}{b_{\min}},\qquad b_{\max}\approx\frac v\omega .
\end{equation}
The logarithm is the same Coulomb logarithm as in the stopping power. The
classical derivation requires $\eta=Ze^2/\hbar v\gg1$ so that $b_{\min}$ exceeds
the quantum length $\hbar/mv$; in the opposite limit $\eta\ll1$ the perturbative
(Born) calculation, not given here, replaces $b_{\min}$ by $\hbar/mv$. The
classical treatment also requires $\hbar\omega\ll E$, since a classical
spectrum cannot carry more energy than the electron has.

\section{Problems}
\begin{enumerate}
  \item Derive \eqref{eq:impulse} and show that the longitudinal impulse
        vanishes.
  \item Evaluate $\int_{\xi_{\min}}^\infty\xi(K_1^2+K_0^2)\,\dd\xi$ numerically for
        $\xi_{\min}=10^{-3}$ and compare with $\ln(1.123/\xi_{\min})$.
  \item Derive \eqref{eq:sigcap} from the condition $V_m=E$, and verify the cases
        $n=4$ and $n\to2$.
  \item For a proton with $v=10^{-2}c$ passing through helium, find $\kB$ and
        decide whether Bohr's or Bethe's formula is appropriate.
\end{enumerate}

\section*{Notes and references}
Bohr's classical theory of energy loss \cite{Bohr13,Bohr48} is derived in \cite{Jackson} (Chapter 13), including the Bessel-function energy transfer; Bethe's quantum formula \cite{Bethe30} and Bloch's correction \cite{Bloch33} are discussed in Chapter~5. Thomson's classical ionization cross section is \cite{Thomson12}. Langevin's capture formula is \cite{Langevin05}; its application to ion--molecule reactions is in \cite{Gioumousis58}, see also \cite{LevineBernstein}. Classical bremsstrahlung goes back to Kramers \cite{Kramers23}; the soft-photon theorem is discussed in \cite{Jackson} (Chapter 15).

% =====================================================================
\part{Quantum Inelastic Scattering}
% =====================================================================

\chapter{Multichannel Scattering Theory}

\section{Channels and coupled equations}

Let the target have internal states $\phi_n(\xi)$ of energy $\varepsilon_n$ and
expand the wave function as $\Psi=\sum_nF_n(\vr)\phi_n(\xi)$, where $\vr$ is the
relative coordinate. With $V_{nm}(\vr)=\langle\phi_n|V|\phi_m\rangle$ the
Schr\"odinger equation becomes a set of coupled equations,
\begin{equation}\label{eq:cc}
  \Bigl[-\frac{\hbar^2}{2\mu}\nabla^2+\varepsilon_n-E\Bigr]F_n(\vr)
  +\sum_mV_{nm}(\vr)F_m(\vr)=0,
\end{equation}
exact for a complete set (this is the expansion, not a perturbative truncation).
The channel wave number is $k_n^2=2\mu(E-\varepsilon_n)/\hbar^2$; the channel is
\emph{open} for $k_n^2>0$ and \emph{closed} otherwise. For a central coupling,
projecting on partial waves gives
\begin{equation}\label{eq:cc-radial}
  \Bigl[\frac{\dd^2}{\dd r^2}-\frac{\ell(\ell+1)}{r^2}+k_n^2\Bigr]u_n
  =\frac{2\mu}{\hbar^2}\sum_mV_{nm}(r)\,u_m(r).
\end{equation}
The long-range character of the problem enters through the asymptotics of the
couplings: a net charge gives a diagonal Coulomb term $\alpha/r$ in the channels
concerned, a charge interacting with a target multipole of order $\lambda$ gives an
off-diagonal tail $V_{nm}\sim r^{-\lambda-1}$, and, most importantly, a charge
interacting with the dipole of a target that has degenerate states of opposite
parity gives $V_{nm}\sim r^{-2}$, an inverse-square coupling that survives
diagonalization (Section~\ref{sec:GD}). Asymptotically,
$F_n\to\delta_{n0}\ee^{\ii k_0z}+f_{n0}\ee^{\ii k_nr}/r$ for short-range
couplings, with the Coulomb-distorted form of the companion volume on elastic scattering
when charges are involved, and the cross sections are given by \eqref{eq:sigfi}.

\section{Unitarity with inelastic channels}\label{sec:unitarity}

Probability conservation makes the full S-matrix, with all channels included,
unitary: $S^\dagger S=1$. Consider the elastic channel in a partial wave $\ell$
and write $S_\ell=\varrho_\ell\,\ee^{2\ii\delta_\ell}$ with $0\le\varrho_\ell\le1$
and a real $\delta_\ell$; $\varrho_\ell<1$ means that flux is lost to other
channels. Then
\begin{align}
  \sigma_{\mathrm{el}}&=\frac{\pi}{k^2}\sum_\ell(2\ell+1)\,|1-S_\ell|^2,\label{eq:pw-el}\\
  \sigma_{\mathrm{inel}}&=\frac{\pi}{k^2}\sum_\ell(2\ell+1)\,\bigl(1-\varrho_\ell^2\bigr),\label{eq:pw-inel}\\
  \sigma_{\mathrm{tot}}&=\frac{2\pi}{k^2}\sum_\ell(2\ell+1)\,\bigl(1-\Real S_\ell\bigr)
  =\sigma_{\mathrm{el}}+\sigma_{\mathrm{inel}}.\label{eq:pw-tot}
\end{align}
Equation \eqref{eq:pw-tot} is the optical theorem in partial waves. The
inelastic formula \eqref{eq:pw-inel} depends only on $|S_\ell|$ and therefore
retains its form when Coulomb waves replace plane waves, even though the
elastic cross section is then infinite. Three consequences follow.
\begin{itemize}
  \item \textbf{Unitarity bounds.} Each partial wave obeys
        $0\le\sigma_{\mathrm{inel}}^{(\ell)}\le\pi(2\ell+1)/k^2$ and
        $\sigma_{\mathrm{el}}^{(\ell)}\le4\pi(2\ell+1)/k^2$. The inelastic cross
        section is finite only if $1-\varrho_\ell^2$ falls faster than $1/\ell$.
  \item \textbf{Absorption implies elastic scattering.} Since $|1-S_\ell|\ge1-\varrho_\ell$,
        a partial wave cannot absorb without also scattering elastically.
  \item \textbf{The black disk.} If $S_\ell=0$ for $\ell\le ka$ and $S_\ell=1$
        beyond, then $\sigma_{\mathrm{inel}}=\pi a^2$ and, from \eqref{eq:pw-el},
        $\sigma_{\mathrm{el}}=\pi a^2$ also; $\sigma_{\mathrm{tot}}=2\pi a^2=2\sigma_{\mathrm{inel}}$.
        The elastic part is shadow (diffraction) scattering, required for the
        wave to be absent behind the absorber. A \emph{quantum} absorber with a
        classical capture region of cross section $\sigma_{\mathrm{cap}}$ is
        therefore accompanied by an equal elastic cross section; the classical
        formulas of Section~\ref{sec:capture} give the inelastic part only.
\end{itemize}

\section{Reciprocity and detailed balance}

Time-reversal invariance makes the S-matrix, in a suitable basis of states of
definite total angular momentum, symmetric: $S_{fi}=S_{if}$. Summing and averaging
over the degenerate states of each channel, with statistical weights $g_i$, $g_f$
(the numbers of internal and spin states), gives the exact relation
\begin{equation}\label{eq:detailed}
  g_i\,k_i^2\,\sigma_{i\to f}=g_f\,k_f^2\,\sigma_{f\to i},
\end{equation}
at the same total energy. It connects excitation and de-excitation (``superelastic'')
cross sections. At the threshold of an exoergic reaction $i\to f$ with $k_i\to0$,
\eqref{eq:detailed} shows that $\sigma_{i\to f}\propto1/k_i\propto1/v$, the
familiar $1/v$ law, provided the reverse process $f\to i$ has a finite cross
section there.

\section{Feshbach projection and the optical potential}

To focus on the elastic channel, split the space with projectors $P$ onto the
channel(s) of interest and $Q=1-P$ onto the rest. The Schr\"odinger equation
$(E-H)\Psi=0$ becomes
$(E-H_{PP})\Psi_P=H_{PQ}\Psi_Q$ and $(E-H_{QQ})\Psi_Q=H_{QP}\Psi_P$. Solving the
second with outgoing-wave boundary conditions and substituting in the first gives
an \emph{exact} one-channel equation,
\begin{equation}\label{eq:optical-pot}
  \bigl[E-T-U_{\mathrm{opt}}(E)\bigr]\Psi_P=0,\qquad
  U_{\mathrm{opt}}=V_{PP}+V_{PQ}\,\frac{1}{E-H_{QQ}+\ii0}\,V_{QP}.
\end{equation}
The optical potential is non-local and energy-dependent, and above the lowest
$Q$-space threshold it is complex, because
$\Imag(E-H_{QQ}+\ii0)^{-1}=-\pi\delta(E-H_{QQ})$ makes
$\Imag U_{\mathrm{opt}}=-\pi V_{PQ}\,\delta(E-H_{QQ})\,V_{QP}\le0$. A negative
imaginary part is absorption of flux from the elastic channel. The continuity
equation $\nabla\cdot\vec j=(2/\hbar)\Imag U\,|\Psi|^2$ gives, for a local
$U_{\mathrm{opt}}$ (the non-local case replaces the integrand by
$\Imag\langle\Psi|U|\Psi\rangle$),
\begin{equation}\label{eq:sig-abs}
  \sigma_{\mathrm{inel}}=-\frac{2}{\hbar v}\int\Imag U_{\mathrm{opt}}(\vr)\,
  |\Psi_P^+(\vr)|^2\,\dd^3r .
\end{equation}
Equations \eqref{eq:optical-pot}--\eqref{eq:sig-abs} are exact; all the difficulty
is hidden in $Q$-space. Their value is structural: they show that inelasticity acts
on the elastic channel as a complex potential whose imaginary part is the
absorption.

The \emph{real} part has an exact long-range form. A slow charge $ze$ at distance
$r$ polarizes a neutral target of polarizability $\alpha_d$, inducing a dipole
$\alpha_dze/r^2$ with interaction energy $-\tfrac12\alpha_d(ze/r^2)^2$, so
\begin{equation}\label{eq:Vpol}
  \mathrm{Re}\,U_{\mathrm{opt}}\to-\frac{\alpha_d\,z^2e^2}{2r^4}\qquad(r\to\infty),
\end{equation}
the polarization potential. It is the origin of Langevin capture
(Section~\ref{sec:capture}), valid when the projectile is slow compared with the
internal motion of the target.

\section{Opening of a new channel: the Wigner cusp}

Consider the elastic S-wave amplitude $S_{ii}$ near the threshold of an S-wave
channel $f$ that opens at energy $E_{\mathrm{th}}$, with $k_f\to0$. The matrix
elements are analytic in $k_f$ (not in $E$): $S_{ii}=a+\beta k_f+\dots$ with
$|a|=1$, and $S_{if}=c\sqrt{k_f}+\dots$ (the $\sqrt{k_f}$ is the flux factor in
\eqref{eq:sigfi}). Unitarity for the row $i$ above threshold, with $k_f$ real,
requires $|S_{ii}|^2+|S_{if}|^2=1$, that is $2\Real\beta+|c|^2=0$. Below
threshold, $k_f=\ii\kappa_f$ is imaginary, channel $f$ is closed, only the
elastic channel is open, and $|S_{ii}|=1$ requires $\Imag\beta=0$. Therefore
$\beta=-|c|^2/2$ is real and fixed. Writing $a=\ee^{2\ii\delta^0}$,
\begin{equation}\label{eq:cusp}
  \Delta\sigma_{\mathrm{el}}=
  \begin{cases}
   -2\sin^2\delta^0\,\sigma_{if}, & E>E_{\mathrm{th}},\\[2pt]
   \ -\dfrac{\pi}{k_i^2}\,|c|^2\kappa_f\,\sin2\delta^0, & E<E_{\mathrm{th}},
  \end{cases}
  \qquad \sigma_{if}=\frac{\pi}{k_i^2}|c|^2k_f .
\end{equation}
Above threshold the elastic cross section \emph{drops} by exactly
$2\sin^2\delta^0$ times the newly opened inelastic cross section, with an
infinite slope $\propto\sqrt{E-E_{\mathrm{th}}}$; below threshold it has a
square-root cusp of either sign. This Wigner cusp is a nonperturbative
consequence of unitarity and analyticity. One checks that the total cross section
changes by $(\pi/k_i^2)|c|^2k_f\cos2\delta^0$, consistent with
$\Delta\sigma_{\mathrm{el}}+\sigma_{if}$. Long-range forces in channel $f$ modify
the analytic dependence on $k_f$ (Chapter~4).

\section{Problems}
\begin{enumerate}
  \item Derive \eqref{eq:pw-inel} from the flux lost from the partial wave
        $\ell$, and verify the black-disk result $\sigma_{\mathrm{el}}=\sigma_{\mathrm{inel}}$.
  \item Using \eqref{eq:detailed}, relate the cross section for $2^3S\to1^1S$
        de-excitation of helium by electrons to the cross section for the
        reverse excitation.
  \item Derive \eqref{eq:cusp} in full, including the check on the total cross
        section.
\end{enumerate}

\section*{Notes and references}
Multichannel scattering theory is developed in \cite{Newton,Taylor,GW,Messiah}; R-matrix and close-coupling methods in \cite{Burke}. The projection-operator formalism and the optical potential are due to Feshbach \cite{Feshbach58} (see \cite{Feshbach92} for a textbook account). Threshold cusps were discussed by Wigner \cite{Wigner48} and Baz' \cite{Baz57}. The detailed-balance relation between excitation and de-excitation of atoms by electrons goes back to Klein and Rosseland \cite{KleinRosseland}; see also \cite{LL3,MM}. The polarization potential is discussed in \cite{MM,BJ}.

% =====================================================================
\chapter{Long-Range Forces and Thresholds}
% =====================================================================

\section{Threshold laws with a Coulomb tail}\label{sec:thr-coulomb}

Near the threshold of channel $f$ with orbital angular momentum $\ell_f$, the
transition amplitude is controlled by the value of the final-channel radial
function in the region $r\lesssim r_0$ where the coupling acts, with
$k_fr_0\ll1$. For short-range forces $u_f\propto(k_fr)^{\ell_f+1}$ and, with the
flux normalization, $\sigma_{if}\propto k_f^{2\ell_f+1}$ (Wigner's law; the
$S$-wave case $\propto k_f$ was used in Section~3.5). If channel $f$ contains a
Coulomb tail, the Coulomb function that is normalized to unit amplitude at
infinity has, at small $\rho=k_fr$, the form $F_\ell\simeq C_\ell(\eta_f)\rho^{\ell+1}$ with
\begin{equation}\label{eq:Cl}
  C_\ell(\eta)=\frac{2^\ell\,\ee^{-\pi\eta/2}\,|\Gamma(\ell+1+\ii\eta)|}{(2\ell+1)!},
  \qquad C_0^2=\frac{2\pi\eta}{\ee^{2\pi\eta}-1}.
\end{equation}
Therefore the threshold law becomes
\begin{equation}\label{eq:thr-law}
  \sigma_{if}\ \propto\ k_f^{\,2\ell_f+1}\,C_{\ell_f}^2(\eta_f),\qquad
  \eta_f=\frac{\mu_f\alpha_f}{\hbar^2k_f},
\end{equation}
with three qualitatively different regimes.
\begin{itemize}
  \item \textbf{Neutral final channel} ($\eta_f=0$): Wigner's power law
        $\sigma\propto k_f^{2\ell+1}$, vanishing at threshold.
  \item \textbf{Repulsive final channel} ($\eta_f>0$): $C_\ell^2\sim\ee^{-2\pi\eta_f}$
        and the cross section vanishes faster than any power,
        $\sigma_{if}\propto\exp(-2\pi\eta_f)$, the Gamow factor of tunnelling
        out of the Coulomb barrier (nuclear reactions with charged outgoing
        fragments).
  \item \textbf{Attractive final channel} ($\eta_f<0$, an electron leaving an
        ion): $|\Gamma(\ell+1+\ii\eta)|^2\simeq2\pi|\eta|^{2\ell+1}\ee^{-\pi|\eta|}$ so
        $C_\ell^2\propto|\eta_f|^{2\ell+1}\propto k_f^{-(2\ell+1)}$, which cancels
        the Wigner factor exactly: the cross section tends to a \emph{finite
        constant} at threshold, a step. For $\ell_f=0$ the step is
        \begin{equation}\label{eq:seaton-step}
          \sigma_{if}\to\frac{2\pi^2|c|^2}{k_i^2\,a_f},\qquad
          a_f=\frac{\hbar^2}{\mu_f|\alpha_f|},
        \end{equation}
        where $a_f$ is the Bohr radius of the final channel and $c$ is the
        short-range coupling constant of Section~3.5.
\end{itemize}
The step \eqref{eq:seaton-step} is the Coulomb version of the Wigner cusp: the
inelastic cross section does not vanish at threshold, so the elastic cross
section drops abruptly there instead of with a square-root singularity. It is
why electron--ion excitation cross sections are finite at threshold while
electron--atom ones vanish.

\section{Below threshold: quantum defects and Rydberg series}\label{sec:qdt}

The attractive Coulomb tail has infinitely many bound states, and the closed
channels below a threshold form Rydberg series that converge to it. Outside the
core of an atom or ion, at $r>r_c$, the radial function is a combination of the
regular and irregular Coulomb functions $f_\ell$ and $g_\ell$ of the pure
$-e^2/r$ field. Since the core region has energy scales of order atomic units,
while the Rydberg series spans energies of order $\mathrm{Ry}/n^2$, the
matching to the core region is the same for all the Rydberg levels, and can be
summarised by a slowly varying \emph{quantum defect} $\mu_\ell(E)$:
\begin{equation}\label{eq:qdt}
  u_\ell\propto f_\ell(E,r)\cos\pi\mu_\ell+g_\ell(E,r)\sin\pi\mu_\ell .
\end{equation}
For $E<0$, requiring decay at infinity gives the levels
\begin{equation}\label{eq:qd-levels}
  E_{n\ell}=-\frac{\mathrm{Ry}}{(n-\mu_\ell)^2},
\end{equation}
and for $E>0$ the same function gives a scattering phase shift, measured against
the pure Coulomb waves, of $\delta_\ell=\pi\mu_\ell$. Since $\mu_\ell(E)$ is smooth
through $E=0$, this is \emph{Seaton's theorem}: the phase shift just above the
ionization threshold equals $\pi$ times the quantum defect of the Rydberg series
just below it, and the bound and continuum spectra are two aspects of a single
function. With coupled channels the quantum defects become a matrix and the
Rydberg series of a closed channel become resonances in the open channels.
Because the Rydberg wave function in the core region has density $\propto\nu^{-3}$
($\nu=n-\mu$), the resonance widths scale as
\begin{equation}\label{eq:gamma-nu3}
  \Gamma_n=\frac{\Gamma^*}{\nu_n^3},\qquad
  \frac{\Gamma_n}{|E_{n+1}-E_n|}\to\frac{\Gamma^*}{2\,\mathrm{Ry}},
\end{equation}
since the level spacing is also $\propto\nu^{-3}$. The ratio of width to spacing
is independent of $n$, so the resonances overlap into a smooth continuum at
threshold, and the cross section averaged over the resonances continues smoothly
across it. Each resonance has the asymmetric shape found by Fano from the exact
diagonalization of one discrete state coupled to a continuum:
\begin{equation}\label{eq:fano}
  \sigma(E)=\sigma_a\,\frac{(q+\epsilon)^2}{1+\epsilon^2}+\sigma_b,\qquad
  \epsilon=\frac{2(E-E_r)}{\Gamma},
\end{equation}
with the asymmetry parameter $q$ and the width $\Gamma$ determined by the
coupling. This is the long-range Coulomb origin of the dense resonance
structure in electron--ion collisions.

\section{The Wannier threshold law for ionization}\label{sec:wannier}

Ionization of an atom or ion by electron impact leaves three charged particles
interacting through long-range Coulomb forces: two electrons and the residual ion
of charge $Z$. Section~2.3 gave Thomson's classical cross section, which rises
linearly with the excess energy $\varepsilon=E-I$, and phase space for three free
bodies would give $\varepsilon^2$ (the non-relativistic $n$-body phase space
scales as $\varepsilon^{(3n-5)/2}$). Wannier (1953) showed that the Coulomb forces
change the exponent.

Take the nucleus fixed, let the electrons have distances $r_1,r_2$ from it and mutual
angle $\theta_{12}$, and define the hyperradius $R=\sqrt{r_1^2+r_2^2}$ and the
hyperangle $\tan\alpha=r_2/r_1$. The potential energy is
$V=-e^2\mathcal C(\alpha,\theta_{12})/R$ with
\begin{equation}\label{eq:Wannier-C}
  \mathcal C(\alpha,\theta_{12})=\frac{Z}{\cos\alpha}+\frac{Z}{\sin\alpha}
  -\frac1{\sqrt{1-\sin2\alpha\cos\theta_{12}}} .
\end{equation}
At the \emph{Wannier ridge} $r_1=r_2$, $\theta_{12}=\pi$ (the two electrons on
opposite sides of the nucleus at equal distance), $\mathcal C=(4Z-1)/\sqrt2$. The
ridge is a saddle point of $\mathcal C$: expanding to second order shows it is a
minimum of the potential in $\theta_{12}$ (stable) and a maximum in $\alpha$
(unstable: if one electron moves ahead it takes more of the energy and the other
falls back). With $R$ playing the role of time, the deviations from the ridge
grow or oscillate as powers of $R$, and the trajectories that lead to
double escape are those that stay near the ridge through the Coulomb zone
$R\lesssim e^2/\varepsilon$. The fraction of phase space that does so scales as a
power of $\varepsilon$ determined by the instability exponent, with the result
\begin{equation}\label{eq:wannier}
  \sigma_{\mathrm{ion}}\propto\varepsilon^{\lambda},\qquad
  \lambda=\frac14\left[\frac{100Z-9}{4Z-1}\right]^{1/2}-\frac14 .
\end{equation}
For electron impact on hydrogen ($Z=1$), $\lambda=1.127$; for ionization of
He$^+$ ($Z=2$), $\lambda=1.056$. The exponent lies between the Thomson value $1$
and the free-phase-space value $2$, and it is a property of the $1/r$ forces of all
three pairs. It is a classical result; quantum calculations confirm the same
exponent very close to threshold. We quote \eqref{eq:wannier} rather than
re-derive the linear stability analysis here.

\section{Dipole coupling and a geometric series of resonances}\label{sec:GD}

When degenerate target states of opposite parity are coupled by a charge--dipole
interaction, diagonalization of $V_{nm}\sim r^{-2}$ gives effective
inverse-square potentials $g/r^2$ in the eigenchannels, with $g$ determined by the
dipole matrix element. For the radial equation $u''+[k^2-g_{\mathrm{eff}}/r^2]u=0$
with $g_{\mathrm{eff}}=\ell(\ell+1)+g$, the behaviour at small $r$ changes
character when $g_{\mathrm{eff}}<-1/4$ (the fall-to-centre condition of the
inverse-square potential). Writing $g_{\mathrm{eff}}=-(\tfrac14+s^2)$, the
solutions are
\begin{equation}\label{eq:logperiodic}
  u\propto\sqrt r\,\cos(s\ln r+\varphi_0),
\end{equation}
which oscillates infinitely often as $r\to0$. Matching to the outer region at
$kr\sim1$ requires $s\ln(1/k)=n\pi+\text{const}$, so resonances (or bound states)
occur at
\begin{equation}\label{eq:GD}
  E_n=E_0\,\ee^{-2\pi n/s},\qquad\frac{E_{n+1}}{E_n}=\ee^{-2\pi/s},
\end{equation}
an infinite geometric series accumulating at threshold (Gailitis and Damburg,
1963). The cross sections oscillate with the logarithm of the energy with period
$2\pi/s$. The mechanism is the same scale invariance of the $1/r^2$ potential
that gave the fall to the centre in the elastic problem.

\section{Problems}
\begin{enumerate}
  \item Derive the small-$\rho$ behaviour $F_0\simeq C_0\rho$ and the formula for
        $C_0^2$ from the Coulomb wave function normalization.
  \item Show that for an attractive final channel with $\ell_f=1$ the threshold
        law \eqref{eq:thr-law} again gives a finite step.
  \item Verify that $\mathcal C=(4Z-1)/\sqrt2$ at the ridge, and that the
        second-order expansion is $\mathcal C_0+\frac{12Z-1}{2\sqrt2}a^2-\frac{1}{8\sqrt2}\phi^2$
        for $\alpha=\pi/4+a$, $\theta_{12}=\pi-\phi$.
  \item Compute the ratio $E_{n+1}/E_n$ in \eqref{eq:GD} for $s=0.5$.
\end{enumerate}

\section*{Notes and references}
Threshold laws are due to Wigner \cite{Wigner48}; their modification by Coulomb and dipole potentials is reviewed in \cite{Sadeghpour00}. Quantum defect theory is due to Seaton \cite{Seaton58,Seaton83}; see also \cite{Fano70,GFS79,GRF82} and the monographs \cite{FanoRau,Gallagher}. The asymmetric resonance profile is due to Fano \cite{Fano61}. Wannier's classical threshold law is \cite{Wannier53}; a quantum-mechanical derivation of the same exponent was given by Rau \cite{Rau71}. The dipole-induced oscillations are due to Gailitis and Damburg \cite{GD63}; the close-coupling treatment is in \cite{Burke}.

% =====================================================================
\chapter{The Semiclassical Impact-Parameter Method}
% =====================================================================

\section{Classical trajectory, quantum internal state}

When the relative motion is classical ($\eta\gg1$) and the energy transferred is
small compared with the collision energy ($\Delta E\ll E$), the heavy particles
follow a classical trajectory $\vr(t)$, unchanged by the excitation, while the
internal state evolves quantum mechanically in the time-dependent field
$V(\vr(t))$. Expanding in the internal states $\phi_n$ with
$\omega_{nm}=(\varepsilon_n-\varepsilon_m)/\hbar$,
\begin{equation}\label{eq:sc-eqs}
  \ii\hbar\,\dot c_n=\sum_mV_{nm}(\vr(t))\,\ee^{\ii\omega_{nm}t}\,c_m ,
\end{equation}
which is exact for given $\vr(t)$ and conserves $\sum_n|c_n|^2=1$. The
probability of reaching state $n$ in a collision with impact parameter $b$ is
$P_n(b)=|c_n(\infty)|^2$ and
\begin{equation}\label{eq:sc-xs}
  \sigma_n=\int2\pi b\,P_n(b)\,\dd b\quad\text{(straight line),}\qquad
  \sigma_n=\int\dd\sigma_{\mathrm{R}}(\theta)\,P_n(\theta)\quad\text{(Rutherford orbit).}
\end{equation}
Equation \eqref{eq:sc-eqs} is non-perturbative in the coupling: it can be solved
exactly or numerically for strong fields. The method needs $\eta\gg1$ so that the
trajectory is classical; it is the quantum counterpart of the classical energy
transfer of Chapter~2 in which the target is also quantum.

\section{The forced harmonic oscillator: an exact result}\label{sec:forced}

Let the target be a harmonic oscillator of frequency $\omega$ driven by a force
$F(t)$: $H=\hbar\omega(a^\dagger a+\tfrac12)-xF(t)$. The solution is exact: starting
in the ground state, the oscillator ends in a \emph{coherent state}, and the
probability of finding it in level $n$ after the force has acted is Poissonian,
\begin{equation}\label{eq:poisson}
  P_n=\frac{\ee^{-\lambda}\lambda^n}{n!},\qquad
  \lambda=\frac{|\tilde F(\omega)|^2}{2m\hbar\omega}=\frac{\Delta E_{\mathrm{cl}}}{\hbar\omega},
\end{equation}
where $\Delta E_{\mathrm{cl}}$ is the classical energy transfer. The mean energy
transferred is therefore exactly the classical \eqref{eq:DEosc}, for every
strength of coupling. For a passing charge on a straight line,
\begin{equation}\label{eq:lambda-b}
  \lambda(b)=\frac{b_0^2}{b^2}\,\xi^2\bigl[K_1^2(\xi)+K_0^2(\xi)\bigr],\qquad
  b_0^2=\frac{2z^2e^4}{m\hbar\omega v^2},\quad\xi=\frac{\omega b}{v}.
\end{equation}
Two regimes appear. For $b\ll v/\omega$, $\lambda\simeq b_0^2/b^2$: the excitation
probability $1-\ee^{-\lambda}$ saturates at 1 for $b\lesssim b_0$ and falls as
$b_0^2/b^2$ beyond. For $b\gg v/\omega$, $\lambda\propto\ee^{-2\omega b/v}$: the
adiabatic regime, with exponentially small probabilities.

\section{Finite inelastic cross sections}

The total inelastic cross section of the oscillator is
\begin{equation}\label{eq:sig-osc}
  \sigma_{\mathrm{inel}}=\int2\pi b\,\bigl[1-\ee^{-\lambda(b)}\bigr]\dd b
  \simeq\pi b_0^2\Bigl[2\ln\frac{v}{\omega b_0}+O(1)\Bigr],
\end{equation}
with $\sigma_n=\int2\pi b\,P_n\,\dd b$ for individual levels. The integral
converges at both ends: at small $b$ the quantum saturation $P\le1$ cuts off the
$1/b^2$ growth of the classical energy transfer at $b\sim b_0$, and at large $b$
the adiabatic factor cuts off the $1/b^2$ tail at $b\sim v/\omega$. This is the
semiclassical realization of the general statement of Section~1.4: the
inelastic cross section is finite for an infinite-range force, while the
elastic Coulomb cross section is not. Higher excitations $n\gg1$ occur for
$b\ll b_0$, with the characteristic Poisson width.

\section{Bohr, Bethe and Bloch: the role of $b_{\min}$}\label{sec:bloch}

Because the mean energy transfer of an oscillator is exactly classical
(Section~\ref{sec:forced}), the quantum mechanics of the target does not change
the stopping power \eqref{eq:Bohr-int}. What differs between the classical and
quantum treatments is the lower cutoff. Classically $b_{\min}=ze^2/mv^2$ is set by
the break-down of the impulse approximation; quantum mechanically the projectile
cannot transfer momentum larger than $2mv$, so the smallest meaningful impact
parameter is $\sim\hbar/mv$. The effective cutoff is the larger of the two, and
their ratio is $b_{\min}^{\mathrm{cl}}/b_{\min}^{\mathrm{q}}=\kB/2$. Hence:
\begin{itemize}
  \item for $\kB\gg1$ the classical value controls and one finds Bohr's formula
        \eqref{eq:Bohr}, whose logarithm contains the constant $1.123$;
  \item for $\kB\ll1$ the quantum value controls and one finds Bethe's formula
        $\ln(2mv^2/I)$ (quoted, since it is a Born-approximation result);
  \item at intermediate $\kB$ Bloch's formula (quoted) interpolates between them,
        \begin{equation}\label{eq:bloch}
          L=\ln\frac{2mv^2}{I}+\psi(1)-\Real\,\psi\!\Bigl(1+\ii\,\frac{\kB}{2}\Bigr),
        \end{equation}
        where $\psi$ is the digamma function and $S=(4\pi z^2e^4NZ/mv^2)L$.
\end{itemize}
The two limits of \eqref{eq:bloch} can be checked: for $\kB\to0$,
$\Real\,\psi(1+\ii y)-\psi(1)=O(y^2)\to0$ and $L\to\ln(2mv^2/I)$, while for
$\kB\to\infty$, $\Real\,\psi(1+\ii y)\to\ln y$ and
$L\to\ln\bigl[2mv^3\hbar/(Ize^2)\bigr]-\gamma_E=\ln\bigl[1.123\,mv^3/(ze^2\omega)\bigr]$
with $I=\hbar\omega$, which is Bohr's formula \eqref{eq:Bohr} including the
constant $1.123$. The digamma function is the same one that appeared in the exact
elastic Coulomb problem; the parameter $\kB=2\eta$ governs the classical limit in
both.

\section{Nonadiabatic transitions: the Landau--Zener model}\label{sec:LZ}

Slow collisions are better described in a basis of states that follow the
instantaneous interaction. When two \emph{diabatic} potential curves cross at a
separation $R_x$, with diabatic energies $H_{11}(R)$, $H_{22}(R)$ and a coupling
$H_{12}$, the system can either stay on the diabatic curve (jumping between
adiabatic curves) or follow the adiabatic one. Near the crossing the diabatic
energy difference changes linearly in time, $\varepsilon(t)=\dot\varepsilon\,t$
with $\dot\varepsilon=v_r|\Delta F|$, where $v_r$ is the radial velocity at the
crossing and $\Delta F=\dd(H_{11}-H_{22})/\dd R$. The Landau--Zener model
(constant $H_{12}$, linear $\varepsilon$) is exactly solvable and gives the
probability of remaining diabatic in one passage,
\begin{equation}\label{eq:LZ}
  P_D=\exp\Bigl(-\frac{2\pi H_{12}^2}{\hbar\,v_r\,|\Delta F|}\Bigr).
\end{equation}
The result is non-perturbative in $H_{12}$: for strong coupling or slow motion
$P_D\to0$ (adiabatic following), for weak coupling or fast motion $P_D\to1$. A
collision passes the crossing twice, on the way in and on the way out, so the
probability of ending on the other adiabatic curve is
\begin{equation}\label{eq:LZ2}
  P_{12}=2P_D(1-P_D)\qquad\text{(phase-averaged)},
\end{equation}
with Stueckelberg interference oscillations $4P_D(1-P_D)\sin^2\Phi$ when the phase
$\Phi$ accumulated between the two passages is retained. The cross section is
obtained by integrating over $b$ with $v_r(b)=v\sqrt{1-U-b^2/R_x^2}$, where
$U=V(R_x)/E$:
\begin{equation}\label{eq:sigLZ}
  \sigma_{12}=4\pi R_x^2\int_0^{\sqrt{1-U}}x\,P_D(x)\,[1-P_D(x)]\,\dd x,\qquad
  P_D(x)=\exp\!\Bigl(-\frac{A}{\sqrt{1-U-x^2}}\Bigr),
\end{equation}
with $x=b/R_x$ and $A=2\pi H_{12}^2/(\hbar v|\Delta F|)$.

\section{Harpooning and the adiabaticity of Coulomb excitation}

Long-range Coulomb forces produce crossings at large separations. In the
collision of an alkali atom A with a halogen B, the covalent (neutral) curve is
nearly flat at large $R$, while the ionic curve ($\mathrm A^+\mathrm B^-$) lies above
it by $\Delta E_\infty=I_{\mathrm A}-\mathcal A_{\mathrm B}$ (ionization energy minus
electron affinity) and falls as the Coulomb attraction:
$V_{\mathrm{ion}}(R)=\Delta E_\infty-e^2/R$. The curves cross at
\begin{equation}\label{eq:harpoon}
  R_x=\frac{e^2}{I_{\mathrm A}-\mathcal A_{\mathrm B}}=\frac{14.40\ \text{eV\,\AA}}{I_{\mathrm A}-\mathcal A_{\mathrm B}} .
\end{equation}
For Na+Cl, $I-\mathcal A=1.53$~eV and $R_x\approx9$~\AA, so electron transfer
(the ``harpoon'' mechanism) occurs at a distance large compared with the atomic
sizes and gives reactive cross sections of order $\pi R_x^2\sim10^2$~\AA$^2$.
The weak coupling $H_{12}$ at such a distance is the reason why the
Landau--Zener probability \eqref{eq:LZ} governs the outcome.

A second use of the same ideas concerns Coulomb excitation of nuclei at
$\eta\gg1$. The relative motion follows the Rutherford orbit, with distance of
closest approach $a=\alpha/(2E)$ in a head-on collision, and \eqref{eq:sc-eqs}
describes the multiple excitation of the target. Whether a transition with energy
$\Delta E$ is efficiently excited is governed by the Massey parameter
\begin{equation}\label{eq:xi-massey}
  \xi=\eta_f-\eta_i\simeq\eta\,\frac{\Delta E}{2E}=\frac{\Delta E\,a}{\hbar v}:
\end{equation}
excitation is strongly suppressed when $\xi\gg1$ (the collision is adiabatic)
and unsuppressed when $\xi\lesssim1$. This is the Coulomb-orbit form of
\eqref{eq:bmax}.

\section{Problems}
\begin{enumerate}
  \item Show that for the forced oscillator the mean number of quanta
        transferred is $\lambda$ in \eqref{eq:poisson}, and that the variance
        is also $\lambda$.
  \item Evaluate the two limits of \eqref{eq:lambda-b} and obtain
        \eqref{eq:sig-osc} by cutting the integral at $b_0$ and $v/\omega$.
  \item Expand $\Real\,\psi(1+\ii y)$ for small and large $y$ and verify the two
        limits of \eqref{eq:bloch}.
  \item For K+Br$_2$ ($I_{\mathrm K}=4.34$~eV, $\mathcal A_{\mathrm{Br_2}}\approx1.0$~eV)
        estimate $R_x$ and the harpoon cross section.
\end{enumerate}

\section*{Notes and references}
The semiclassical impact-parameter method and its application to Coulomb excitation are reviewed in \cite{Alder56,AlderWinther}; collision-induced transitions between molecular states are treated in \cite{MM,BM,Child,NU}. The Landau--Zener formula is \cite{Landau32,Zener32}, with the interference oscillations of Stueckelberg \cite{Stueckelberg32}. The adiabatic criterion is due to Massey \cite{Massey49}. The forced harmonic oscillator and its coherent-state solution are in \cite{Husimi53,Glauber63}. The stopping-power formulas of Bohr, Bethe and Bloch are \cite{Bohr13,Bethe30,Bloch33}; see \cite{Fano63,Inokuti71} for the quantum theory. The harpoon mechanism is discussed in \cite{LevineBernstein,Child}.

% =====================================================================
\part{Relativistic Inelastic Scattering}
% =====================================================================

\chapter{Classical Relativistic Inelastic Scattering}

\section{Energy transfer by a relativistic charge}

The transverse impulse delivered by a fast charge $ze$ to a particle at rest at
impact parameter $b$ is unchanged by relativity. In the rest frame of the target
the fields of the moving charge are, in Gaussian units,
\begin{equation}\label{eq:rel-fields}
  E_\perp(t)=\frac{\gamma ze\,b}{(b^2+\gamma^2v^2t^2)^{3/2}},\qquad
  E_\parallel(t)=-\frac{\gamma ze\,vt}{(b^2+\gamma^2v^2t^2)^{3/2}},
\end{equation}
and $\int eE_\perp\,\dd t=2ze^2/(bv)$: the factor $\gamma$ in the field is
compensated by the Lorentz-contracted duration of the encounter. Thus the
energy transferred to a free electron is again $W=\Delta p_\perp^2/2m=2z^2e^4/(mv^2b^2)$
and, with $W=-t/2m$ and the invariant Rutherford law
$\dd\sigma/\dd t=4\pi z^2e^4/(\beta^2t^2)$,
\begin{equation}\label{eq:rel-dsdW}
  \frac{\dd\sigma}{\dd W}=\frac{2\pi z^2e^4}{mv^2W^2},\qquad W\ll W_{\max}.
\end{equation}
The largest energy a heavy projectile of mass $M$ can transfer to a free electron
is
\begin{equation}\label{eq:Wmax}
  W_{\max}=\frac{2mc^2\beta^2\gamma^2}{1+2\gamma m/M+(m/M)^2}\ \simeq\ 2mc^2\beta^2\gamma^2,
\end{equation}
so the classical minimum impact parameter becomes
$b_{\min}=ze^2/(\gamma mv^2)$, found by equating $W(b)=W_{\max}$. The quantum
cutoff $\hbar/\gamma mv$ has the same ratio $\kB/2$ to $b_{\min}$ as in the
non-relativistic case, because $\gamma$ cancels.

\section{A bound target in a relativistic field}

For an oscillator of frequency $\omega$, the exact energy transfer is again
$|\tilde F(\omega)|^2/2m$. Fourier transforming \eqref{eq:rel-fields}, with
$\xi=\omega b/(\gamma v)$ (the adiabatic parameter now involves the
\emph{contracted} duration $b/\gamma v$):
\begin{equation}\label{eq:rel-DE}
  \Delta E(b)=\frac{2z^2e^4}{mv^2b^2}\,\xi^2\Bigl[K_1^2(\xi)+\frac1{\gamma^2}K_0^2(\xi)\Bigr].
\end{equation}
The longitudinal part is suppressed by $1/\gamma^2$. The adiabatic distance is
$b_{\max}=\gamma v/\omega$ (cf.\ \eqref{eq:bmax}): relativistic charges transfer
energy out to larger distances because their transverse field is boosted.

\section{The relativistic Bohr formula and the relativistic rise}

Integrating \eqref{eq:rel-DE} over impact parameters as in \eqref{eq:Bohr-int}
requires $\int_{\xi_{\min}}^\infty\xi K_1^2\,\dd\xi=\ln(1.123/\xi_{\min})-\tfrac12$
and $\int_0^\infty\xi K_0^2\,\dd\xi=\tfrac12$, so the bracketed integral becomes
$\ln(1.123/\xi_{\min})-\tfrac12+1/(2\gamma^2)=\ln(1.123/\xi_{\min})-\beta^2/2$. With
$\xi_{\min}=\omega b_{\min}/(\gamma v)=\omega ze^2/(\gamma^2mv^3)$,
\begin{equation}\label{eq:rel-Bohr}
  S=\frac{4\pi z^2e^4NZ}{mv^2}\Bigl[\ln\frac{1.123\,\gamma^2mv^3}{z\,e^2\,\omega}
  -\frac{\beta^2}{2}\Bigr].
\end{equation}
Two features distinguish this from the non-relativistic result \eqref{eq:Bohr}.
First, the logarithm contains $\gamma^2$: after its minimum, at $\beta\gamma$ of a few, the stopping power rises logarithmically with energy, the
\emph{relativistic rise}, because the transverse field reaches farther. Second,
the $-\beta^2/2$ term. The rise does not continue indefinitely in a medium: when
$b_{\max}=\gamma v/\omega$ exceeds the interatomic distance the field is screened
by the polarization of the medium (the density effect of Fermi), and the
stopping power saturates at a plateau; we quote this effect without derivation.

\section{Relativistic soft radiation}\label{sec:rel-soft}

For a relativistic charge the radiated spectrum in the low-frequency limit
depends again only on the initial and final velocities, now through
\begin{equation}\label{eq:rel-soft}
  \frac{\dd^2I}{\dd\omega\,\dd\Omega}\bigg|_{\omega\to0}=\frac{e^2}{4\pi^2c}\,
  \left|\frac{\bm n\times(\bm n\times\bm\beta')}{1-\bm n\cdot\bm\beta'}-
        \frac{\bm n\times(\bm n\times\bm\beta)}{1-\bm n\cdot\bm\beta}\right|^2 ,
\end{equation}
where $\bm\beta,\bm\beta'$ are the velocities before and after the collision and
$\bm n$ is the direction of observation. This follows from the exact integral of
the Li\'enard--Wiechert radiation amplitude, whose integrand is a total time
derivative. For $\beta\ll1$, integrating over angles with
$\int|\bm n\times(\bm n\times\Delta\bm\beta)|^2\dd\Omega=(8\pi/3)|\Delta\bm\beta|^2$
reproduces \eqref{eq:soft-nr}. For $\gamma\gg1$ the factors $1/(1-\bm n\cdot\bm\beta)$ concentrate the
radiation in a cone of angle $\sim1/\gamma$ about the velocity, and the spectrum
is flat up to frequencies that grow with $\gamma$. Like \eqref{eq:soft-nr}, the
result is exact classically, but it requires $\hbar\omega\ll E$ and, for Coulomb
collisions, a classical orbit ($\eta\gg1$); at low $\eta$ the perturbative
Bethe--Heitler result applies and is not derived here.

\section{Relativistic capture revisited}

In the elastic volume, the classical orbit of a relativistic particle in an
attractive Coulomb field was shown to spiral into the centre when
$L<|\alpha|/c$, with cross section
\begin{equation}\label{eq:rel-cap}
  \sigma_{\mathrm{cap}}=\frac{\pi\alpha^2}{p^2c^2}.
\end{equation}
This is an absorptive (inelastic) cross section produced by relativity alone,
with no non-relativistic counterpart for a $1/r$ potential: compare the barrier
capture of Section~\ref{sec:capture}, which needs $n>2$. Its quantum meaning is
that $L_c/\hbar=\alpha/\hbar c=Z\afs<1$ for $Z<137$, smaller than the smallest
non-zero orbital angular momentum: for light nuclei the capture region is smaller
than a wavelength and the classical process does not exist. For $Z>137$ it does,
and it connects to the supercritical Coulomb problem of Section~\ref{sec:supercrit}.

\section{Problems}
\begin{enumerate}
  \item Verify \eqref{eq:rel-fields} by Lorentz-transforming the Coulomb field,
        and prove $\int eE_\perp\dd t=2ze^2/(bv)$.
  \item Derive \eqref{eq:Wmax} from the kinematics of a heavy particle colliding
        with a free electron at rest.
  \item Estimate, from \eqref{eq:rel-Bohr}, the factor by which $S$ at $\gamma=100$
        exceeds its value at $\gamma=1$ when $\ln(mv^3/ze^2\omega)=10$ at
        $\gamma=1$, ignoring the density effect.
\end{enumerate}

\section*{Notes and references}
Relativistic energy transfer and the relativistic rise are treated in \cite{Jackson} (Chapter 13); the density effect was first calculated by Fermi \cite{Fermi40}. For comprehensive reviews of the stopping power at high energies see \cite{Sigmund,PDG}. The relativistic soft-photon formula and its derivation from the Li\'enard--Wiechert potentials are in \cite{Jackson,LL2}.

% =====================================================================
\chapter{Relativistic Quantum Inelastic Scattering}
% =====================================================================

\section{Generalized unitarity and phase space}

With the invariant amplitude defined by
$S_{fi}=\delta_{fi}+\ii(2\pi)^4\delta^4(P_f-P_i)\mathcal M_{fi}$ (units
$\hbar=c=1$), unitarity of $S$ for \emph{all} channels gives the exact relation
\begin{equation}\label{eq:gen-unit}
  -\ii\bigl(\mathcal M_{fi}-\mathcal M^*_{if}\bigr)=\sum_n\int\dd\Phi_n\,
  \mathcal M^*_{nf}\,\mathcal M_{ni},
\end{equation}
where the sum runs over every intermediate state $n$ and $\dd\Phi_n$ is the
Lorentz-invariant phase space. Setting $f=i$ and forward kinematics,
\begin{equation}\label{eq:rel-optical}
  2\,\Imag\mathcal M_{ii}(s,t=0)=\sum_n\int\dd\Phi_n\,|\mathcal M_{ni}|^2
  =F\,\sigma_{\mathrm{tot}},\qquad F=4p^*\sqrt s,
\end{equation}
the optical theorem with $\sigma_{\mathrm{tot}}$ summed over all channels. The
two-body phase space is $\Phi_2=p_f^*/(4\pi\sqrt s)$, and near threshold the
phase space of $n$ non-relativistic bodies behaves as
$\Phi_n\propto\varepsilon^{(3n-5)/2}$ with $\varepsilon$ the kinetic energy above
threshold. Relations \eqref{eq:gen-unit}--\eqref{eq:rel-optical} hold with no
expansion in the coupling.

\section{Thresholds, cuts and Coulomb final-state factors}

Equation \eqref{eq:gen-unit} implies that $\mathcal M_{ii}(s)$ has a branch point at
each threshold $s_{\mathrm{th}}=(\sum_fm_f)^2$ and a right-hand cut above it. The
non-relativistic Wigner cusp of Section~3.5 is its low-energy form, with
$k_f\to p_f^*$: near an S-wave threshold the elastic cross section drops by
$2\sin^2\delta^0$ times the new inelastic cross section. The Coulomb tails of
Chapter~4 persist relativistically as the \emph{Sommerfeld--Sakharov factor}: for a pair
of final particles with opposite charges and relative velocity $v_{\mathrm{rel}}$
the threshold cross section is multiplied by
\begin{equation}\label{eq:sakharov}
  S=\frac{2\pi|\eta|}{1-\ee^{-2\pi|\eta|}},\qquad|\eta|=\frac{\alpha}{v_{\mathrm{rel}}},
\end{equation}
which makes the pair-production cross section jump to a finite value at threshold
instead of vanishing like the phase-space factor $\beta$. For like-charged final
particles the corresponding factor $2\pi\eta/(\ee^{2\pi\eta}-1)$ is the
Gamow suppression. These factors are the relativistic images of the threshold
laws \eqref{eq:thr-law} and use the same function $C_0^2(\eta)$.

\section{Impact-parameter unitarity and bounds on cross sections}

At high energy and small angles the amplitude can be written in terms of a profile
function in impact-parameter space, $S(b)=\ee^{2\ii\delta(b)}$ with
$\Imag\delta\ge0$. The cross sections follow exactly from the profile
$\Gamma(b)=1-S(b)$:
\begin{equation}\label{eq:b-unit}
  \sigma_{\mathrm{tot}}=2\!\int\!\dd^2b\,\Real\Gamma,\quad
  \sigma_{\mathrm{el}}=\!\int\!\dd^2b\,|\Gamma|^2,\quad
  \sigma_{\mathrm{inel}}=\!\int\!\dd^2b\,\bigl(1-|S|^2\bigr).
\end{equation}
These are the continuum counterparts of \eqref{eq:pw-el}--\eqref{eq:pw-tot}. For a
black disk of radius $R$ ($S=0$ for $b<R$) they give
$\sigma_{\mathrm{el}}=\sigma_{\mathrm{inel}}=\pi R^2$ and
$\sigma_{\mathrm{el}}/\sigma_{\mathrm{tot}}=\tfrac12$.

Unitarity also bounds how fast total cross sections can grow, provided the range
of the force is finite. Suppose the amplitude has its nearest singularity in $t$ at
the two-pion threshold $t_0=4m_\pi^2$. Then in impact-parameter space the
profile decays as $\ee^{-2m_\pi b}$ times a power of $s$, so $\Gamma(b)$ is
negligible beyond
$b_{\max}\simeq\ln(s/s_0)/(2m_\pi)$, with $s_0$ a scale that unitarity alone does
not fix. Since $|\Gamma|\le2$,
\begin{equation}\label{eq:froissart}
  \sigma_{\mathrm{tot}}\ \lesssim\ 4\pi b_{\max}^2=\frac{\pi}{m_\pi^2}\,\ln^2\frac{s}{s_0},
\end{equation}
the Froissart--Martin bound (our argument is heuristic; the rigorous proof uses
analyticity in $\cos\theta$ within the Lehmann--Martin ellipse). The contrast with
the long-range case is instructive: for a massless exchanged particle there is no
mass gap, no $b_{\max}$, and the \emph{elastic} Coulomb cross section is already
infinite. The finiteness of inelastic cross sections discussed in
Section~\ref{sec:qmin} comes from a different mechanism, the adiabatic cutoff.

\section{Pair creation as Landau--Zener tunnelling in the Dirac sea}\label{sec:schwinger}

The Landau--Zener formula \eqref{eq:LZ} gives a non-perturbative derivation of
vacuum pair creation in a uniform electric field $\mathcal E$. A Dirac particle of
momentum $\bm p$ in a uniform field has kinetic momentum $p_z-e\mathcal Et$. For
each $\bm p$, the diabatic states are the positive-energy and negative-energy
branches. At $p_\perp=0$ they have energies $\pm c(p_z-e\mathcal Et)$, which cross
as the momentum sweeps through zero, and the coupling that opens the gap is the
mass term, $H_{12}=mc^2$. The rate of change of the diabatic energy difference is
$|\dot\varepsilon|=2ce\mathcal E$. The Landau--Zener probability of the system
staying on the diabatic curve, which is the probability that a negative-energy
state becomes a positive-energy electron, i.e.\ that a pair is created, is
\begin{equation}\label{eq:LZ-pair}
  P=\exp\Bigl(-\frac{2\pi m^2c^4}{\hbar\cdot2ce\mathcal E}\Bigr)
   =\exp\Bigl(-\frac{\pi m^2c^3}{\hbar e\mathcal E}\Bigr)=\exp\Bigl(-\pi\frac{\mathcal E_c}{\mathcal E}\Bigr),
\end{equation}
with the critical field $\mathcal E_c=m^2c^3/(e\hbar)=1.3\times10^{18}$~V/m. For
nonzero $p_\perp$ the replacement $m^2\to m^2+p_\perp^2/c^2$ gives a Gaussian in
$p_\perp$. Counting modes, with two spin states, a momentum range $\Delta p_z=e\mathcal ET$ and
$\int\dd^2p_\perp\exp(-\pi cp_\perp^2/e\mathcal E\hbar)=e\mathcal E\hbar/c$, the pair
production rate per unit volume and time is
\begin{equation}\label{eq:schwinger}
  w=2\,\frac{e\mathcal E}{(2\pi\hbar)^3}\,\frac{e\mathcal E\hbar}{c}\,\ee^{-\pi\mathcal E_c/\mathcal E}
   =\frac{(e\mathcal E)^2}{4\pi^3\hbar^2c}\,\ee^{-\pi\mathcal E_c/\mathcal E},
\end{equation}
the leading term of Schwinger's result, whose complete form is a sum over
multiple pairs, $\sum_{n\ge1}n^{-2}\exp(-n\pi\mathcal E_c/\mathcal E)$ (quoted).
The same exponent follows from WKB tunnelling through the gap of width
$2mc^2$, which a uniform field spans over a distance $2mc^2/e\mathcal E$:
\begin{equation}
  2\cdot\frac{1}{\hbar c}\int_{-x_0}^{x_0}\!\sqrt{m^2c^4-(e\mathcal Ex)^2}\,\dd x
  =\frac{\pi m^2c^3}{\hbar e\mathcal E},
\end{equation}
which links the result to the Klein step of the elastic volume: for a step of
height $V_0>2mc^2$ and width $w\gg\hbar/mc$ the transmission is exponentially
small, with the Sauter--Schwinger exponent. The process is essentially
non-perturbative in the field: the rate is non-analytic in the coupling $e$ at
$e=0$ and cannot be expanded in a power series in it.

\section{Supercritical Coulomb fields}\label{sec:supercrit}

The Dirac equation for a point nucleus of charge $Ze$ gives the 1s level
$E=mc^2\sqrt{1-\zeta^2}$ (from the Sommerfeld--Dirac formula with
$n=1$, $\kappa_D=-1$, $\zeta=Z\afs$), which vanishes at $\zeta=1$,
$Z=137$, where the point-nucleus problem becomes singular. For a nucleus of finite
radius (uniformly charged sphere, $R_N\approx1.2A^{1/3}$~fm) the potential is
regular at the origin and the level continues to fall, reaching
$E=-mc^2$, the upper edge of the negative-energy continuum, at the critical
charge $Z_{\mathrm{cr}}\approx173$. For $Z>Z_{\mathrm{cr}}$ the 1s state dives
into the negative-energy continuum and becomes a resonance. It decays by
\emph{spontaneous} positron emission: the vacuum is unstable, with two electrons
filling the 1s shell (spin 1/2) and two positrons escaping. This is a
non-perturbative inelastic process, and in the formalism of Fano
(Section~\ref{sec:qdt}), with a discrete state coupled to a continuum, the width
near threshold grows as
\begin{equation}
  \Gamma(Z)\propto(Z-Z_{\mathrm{cr}})^2,
\end{equation}
to leading order, because the coupling matrix element vanishes linearly at
$Z_{\mathrm{cr}}$. Positron scattering from the supercritical nucleus shows a
resonance of Fano shape at the energy of the diving level. The effect has been
sought in collisions of heavy ions (such as U+U, with combined charge $184$) but
the nuclei remain together only $\sim10^{-21}$~s, and the interpretation of the
experiments is not settled.

\section{Problems}
\begin{enumerate}
  \item Derive \eqref{eq:rel-optical} from \eqref{eq:gen-unit} and check that it
        reduces to \eqref{eq:pw-tot} summed over partial waves.
  \item Show that the exponent in \eqref{eq:LZ-pair} equals the WKB tunnelling
        integral across the gap of the Dirac sea.
  \item Compute $\mathcal E_c=m^2c^3/e\hbar$ in V/m and the corresponding laser
        intensity.
  \item Find the 1s energy of a point nucleus for $Z=100,\,130$ and estimate the
        point at which the level reaches $E=0$.
\end{enumerate}

\section*{Notes and references}
The general unitarity relations and the Lorentz-invariant phase space are treated in \cite{Weinberg,BK,PS}; the analytic S-matrix in \cite{Eden,MartinSpearman}. The Froissart bound is \cite{Froissart61}, with the sharpened version of \cite{Martin63}. Pair creation in a constant field was calculated by Sauter \cite{Sauter31}, Heisenberg and Euler \cite{HE36} and Schwinger \cite{Schwinger51}; see \cite{Brezin70} for oscillating fields and \cite{Dunne05,GelisTanji} for reviews. The Landau--Zener description of the same exponent uses \cite{Landau32,Zener32}. Supercritical Coulomb fields are discussed in \cite{Pomeranchuk45,ZP71,Reinhardt77,Greiner}. Sakharov--Sommerfeld factors are \cite{Sakharov48,Sommerfeld31}.

% =====================================================================
\chapter{Synthesis: Inelastic Collisions in Three Settings}
% =====================================================================

\section{The three descriptions compared}

Two dimensionless parameters control the results: the Bohr parameter
$\kB=2\eta=2ze^2/\hbar v$, which says whether the motion is classical, and the
adiabatic parameter $\xi=\omega b/v$ (with $\gamma$ in the relativistic case),
which says whether a target of excitation frequency $\omega$ responds at impact
parameter $b$.

\medskip
\noindent\small
\begin{tabular}{@{}>{\raggedright\arraybackslash}p{2.3cm}>{\raggedright\arraybackslash}p{3.7cm}>{\raggedright\arraybackslash}p{4.2cm}>{\raggedright\arraybackslash}p{4.2cm}@{}}
\hline
 & \textbf{Classical} & \textbf{Quantum} & \textbf{Relativistic}\\
\hline
Energy transfer to a bound target &
  Bessel formula, \eqref{eq:DEosc}, exact for a harmonic target &
  Poisson distribution with the same mean, \eqref{eq:poisson} &
  \eqref{eq:rel-DE} with contracted duration and $1/\gamma^2$ longitudinal term\\[6pt]
Stopping power &
  Bohr, \eqref{eq:Bohr}, $\kB\gg1$ &
  Bethe, $\kB\ll1$; Bloch interpolation, \eqref{eq:bloch} &
  $\ln(1.123\gamma^2\cdots)-\beta^2/2$, \eqref{eq:rel-Bohr}; density effect\\[6pt]
Cutoffs &
  $b_{\min}=ze^2/mv^2$, $b_{\max}=v/\omega$ &
  $b_{\min}\sim\hbar/mv$, $b_{\max}=\hbar v/\Delta E$ &
  $b_{\min}=ze^2/\gamma mv^2$, $b_{\max}=\gamma v/\omega$\\[6pt]
Capture / absorption &
  barrier capture $\propto E^{-2/n}$, \eqref{eq:sigcap}; none for $1/r$ &
  black disk: $\sigma_{\mathrm{el}}=\sigma_{\mathrm{inel}}$, Section~3.2 &
  $\pi\alpha^2/p^2c^2$, \eqref{eq:rel-cap}; classically absent for $Z<137$\\[6pt]
Thresholds &
  Thomson: linear, \eqref{eq:Thomson} &
  Wigner $k_f^{2\ell+1}$; Gamow; Coulomb step, \eqref{eq:thr-law}; Wannier $\varepsilon^{1.127}$ &
  cusps and Sommerfeld--Sakharov factor, \eqref{eq:sakharov}\\[6pt]
Radiation &
  soft theorem, \eqref{eq:soft-nr}, \eqref{eq:brems} &
  needs $\eta\gg1$ for the classical form &
  \eqref{eq:rel-soft}; beaming within $1/\gamma$\\[6pt]
Strong-field effects &
  none (energy loss only) &
  Landau--Zener, \eqref{eq:LZ}; Fano profiles &
  Schwinger pair creation, \eqref{eq:schwinger}; supercritical decay\\
\hline
\end{tabular}
\normalsize

\section{Five faces of inelasticity with long range}

\begin{enumerate}
  \item \textbf{A built-in cutoff.} Because $q_{\min}=\Delta E/\hbar v\neq0$,
        inelastic cross sections do not suffer the forward divergence of elastic
        Coulomb scattering, \eqref{eq:qmin}, \eqref{eq:bmax}.
  \item \textbf{The Coulomb logarithm.} The same $\ln(b_{\max}/b_{\min})$ appears
        in stopping, bremsstrahlung and the momentum-transfer cross section, with
        $b_{\max}$ from adiabaticity and $b_{\min}$ from the classical orbit or the
        quantum length, whichever is larger, \eqref{eq:Bohr}, \eqref{eq:bloch}.
  \item \textbf{Coulomb factors at thresholds.} Gamow suppression, a finite
        Coulomb step, and the Wannier exponent replace the power laws of
        short-range physics, \eqref{eq:thr-law}, \eqref{eq:wannier}.
  \item \textbf{Rydberg series.} Infinitely many closed channels converge on
        every threshold, with widths $\propto\nu^{-3}$ and a smooth continuation
        across it, \eqref{eq:gamma-nu3}.
  \item \textbf{Non-perturbative transitions.} Landau--Zener crossings, forced
        oscillators, Schwinger pair creation and the supercritical Coulomb centre
        are exactly solvable, strong-field processes where perturbation theory in
        the coupling fails, \eqref{eq:poisson}, \eqref{eq:LZ}, \eqref{eq:schwinger}.
\end{enumerate}

\section{What is, and is not, an accident}

That the mean energy transferred to a harmonic target is exactly classical is an
accident of linearity; the Bohr--Bethe crossover happens entirely through the
lower cutoff $b_{\min}$, governed by $\kB=2\eta$, the same parameter that
controls the elastic classical limit. The coincidence of classical and quantum
Rutherford cross sections in the elastic volume and the exactness of the
classical stopping law at large $\kB$ are two sides of the same fact: the $1/r$
force has a classical limit that is controlled by a single number. Relativity
adds $\gamma$ to the cutoffs and a genuine new process (spontaneous pair creation)
that has no classical analogue and no perturbative expansion; the structural
elements of Part~I (thresholds, channel counting, unitarity, the adiabatic cutoff)
survive intact through all three settings.

\section*{Notes and references}
For wide-ranging treatments of atomic and molecular collisions see \cite{MM,BJ,FanoRau,BM,Child,NU,Burke}; for general scattering theory see \cite{Newton,Taylor,GW}.

\backmatter
\renewcommand{\bibname}{References}
\addcontentsline{toc}{chapter}{References}
\begin{thebibliography}{99}
\bibitem{Alder56} K. Alder, A. Bohr, T. Huus, B. Mottelson and A. Winther, Rev. Mod. Phys. \textbf{28}, 432 (1956).
\bibitem{AlderWinther} K. Alder and A. Winther, \emph{Electromagnetic Excitation: Theory of Coulomb Excitation with Heavy Ions}, North-Holland, Amsterdam (1975).
\bibitem{Baz57} A. I. Baz', Zh. Eksp. Teor. Fiz. \textbf{33}, 923 (1957).
\bibitem{Bethe30} H. Bethe, Ann. Phys. (Leipzig) \textbf{397}, 325 (1930).
\bibitem{Bloch33} F. Bloch, Ann. Phys. (Leipzig) \textbf{408}, 285 (1933).
\bibitem{Bohr13} N. Bohr, Phil. Mag. \textbf{25}, 10 (1913).
\bibitem{Bohr48} N. Bohr, Mat.-Fys. Medd. Dan. Vid. Selsk. \textbf{18}, no.~8 (1948).
\bibitem{BM} B. H. Bransden and M. R. C. McDowell, \emph{Charge Exchange and the Theory of Ion--Atom Collisions}, Clarendon Press, Oxford (1992).
\bibitem{BJ} B. H. Bransden and C. J. Joachain, \emph{Physics of Atoms and Molecules}, 2nd ed., Prentice Hall, Harlow (2003).
\bibitem{Brezin70} E. Br\'ezin and C. Itzykson, Phys. Rev. D \textbf{2}, 1191 (1970).
\bibitem{Burke} P. G. Burke, \emph{R-Matrix Theory of Atomic Collisions}, Springer, Berlin (2011).
\bibitem{BK} E. Byckling and K. Kajantie, \emph{Particle Kinematics}, Wiley, London (1973).
\bibitem{Child} M. S. Child, \emph{Molecular Collision Theory}, Academic Press, London (1974); Dover reprint (1996).
\bibitem{Dunne05} G. V. Dunne, Heisenberg--Euler effective Lagrangians: basics and extensions, in \emph{From Fields to Strings: Circumnavigating Theoretical Physics}, World Scientific, Singapore (2005); arXiv:hep-th/0406216.
\bibitem{Eden} R. J. Eden, P. V. Landshoff, D. I. Olive and J. C. Polkinghorne, \emph{The Analytic S-Matrix}, Cambridge University Press (1966).
\bibitem{Fano61} U. Fano, Phys. Rev. \textbf{124}, 1866 (1961).
\bibitem{Fano63} U. Fano, Annu. Rev. Nucl. Sci. \textbf{13}, 1 (1963).
\bibitem{Fano70} U. Fano, Phys. Rev. A \textbf{2}, 353 (1970).
\bibitem{FanoRau} U. Fano and A. R. P. Rau, \emph{Atomic Collisions and Spectra}, Academic Press, Orlando (1986).
\bibitem{Fermi40} E. Fermi, Phys. Rev. \textbf{57}, 485 (1940).
\bibitem{Feshbach58} H. Feshbach, Ann. Phys. (N.Y.) \textbf{5}, 357 (1958); \textbf{19}, 287 (1962).
\bibitem{Feshbach92} H. Feshbach, \emph{Theoretical Nuclear Physics: Nuclear Reactions}, Wiley, New York (1992).
\bibitem{Froissart61} M. Froissart, Phys. Rev. \textbf{123}, 1053 (1961).
\bibitem{GD63} M. Gailitis and R. Damburg, Proc. Phys. Soc. \textbf{82}, 192 (1963); Zh. Eksp. Teor. Fiz. \textbf{44}, 1644 (1963) [Sov. Phys. JETP \textbf{17}, 1107 (1963)].
\bibitem{Gallagher} T. F. Gallagher, \emph{Rydberg Atoms}, Cambridge University Press (1994).
\bibitem{GelisTanji} F. Gelis and N. Tanji, Prog. Part. Nucl. Phys. \textbf{87}, 1 (2016).
\bibitem{Gioumousis58} G. Gioumousis and D. P. Stevenson, J. Chem. Phys. \textbf{29}, 294 (1958).
\bibitem{Glauber63} R. J. Glauber, Phys. Rev. \textbf{131}, 2766 (1963).
\bibitem{GW} M. L. Goldberger and K. M. Watson, \emph{Collision Theory}, Wiley, New York (1964); Dover reprint (2004).
\bibitem{GFS79} C. H. Greene, U. Fano and G. Strinati, Phys. Rev. A \textbf{19}, 1485 (1979).
\bibitem{GRF82} C. H. Greene, A. R. P. Rau and U. Fano, Phys. Rev. A \textbf{26}, 2441 (1982).
\bibitem{Greiner} W. Greiner, B. M\"uller and J. Rafelski, \emph{Quantum Electrodynamics of Strong Fields}, Springer, Berlin (1985).
\bibitem{Hagedorn} R. Hagedorn, \emph{Relativistic Kinematics}, Benjamin, New York (1963).
\bibitem{HE36} W. Heisenberg and H. Euler, Z. Phys. \textbf{98}, 714 (1936).
\bibitem{Husimi53} K. Husimi, Prog. Theor. Phys. \textbf{9}, 381 (1953).
\bibitem{Inokuti71} M. Inokuti, Rev. Mod. Phys. \textbf{43}, 297 (1971).
\bibitem{Jackson} J. D. Jackson, \emph{Classical Electrodynamics}, 3rd ed., Wiley, New York (1999). Chapters 13 (collisions and energy loss), 14 (radiation by moving charges) and 15 (bremsstrahlung) are the most relevant.
\bibitem{KleinRosseland} O. Klein and S. Rosseland, Z. Phys. \textbf{4}, 46 (1921).
\bibitem{Kramers23} H. A. Kramers, Phil. Mag. \textbf{46}, 836 (1923).
\bibitem{Landau32} L. D. Landau, Phys. Z. Sowjetunion \textbf{2}, 46 (1932).
\bibitem{LL2} L. D. Landau and E. M. Lifshitz, \emph{The Classical Theory of Fields}, 4th ed., Pergamon, Oxford (1975).
\bibitem{LL3} L. D. Landau and E. M. Lifshitz, \emph{Quantum Mechanics: Non-Relativistic Theory}, 3rd ed., Pergamon, Oxford (1977). Chapters XVII and XVIII treat elastic and inelastic collisions.
\bibitem{Langevin05} P. Langevin, Ann. Chim. Phys. \textbf{5}, 245 (1905).
\bibitem{LevineBernstein} R. D. Levine and R. B. Bernstein, \emph{Molecular Reaction Dynamics and Chemical Reactivity}, Oxford University Press (1987).
\bibitem{Martin63} A. Martin, Phys. Rev. \textbf{129}, 1432 (1963); Nuovo Cimento A \textbf{42}, 930 (1966).
\bibitem{MartinSpearman} A. D. Martin and T. D. Spearman, \emph{Elementary Particle Theory}, North-Holland, Amsterdam (1970).
\bibitem{Massey49} H. S. W. Massey, Rep. Prog. Phys. \textbf{12}, 248 (1949).
\bibitem{Messiah} A. Messiah, \emph{Quantum Mechanics}, vols.~1--2, North-Holland, Amsterdam (1961--1962).
\bibitem{MM} N. F. Mott and H. S. W. Massey, \emph{The Theory of Atomic Collisions}, 3rd ed., Clarendon Press, Oxford (1965).
\bibitem{Newton} R. G. Newton, \emph{Scattering Theory of Waves and Particles}, 2nd ed., Springer, New York (1982); Dover reprint (2002).
\bibitem{NU} E. E. Nikitin and S. Ya. Umanskii, \emph{Theory of Slow Atomic Collisions}, Springer, Berlin (1984).
\bibitem{PS} M. E. Peskin and D. V. Schroeder, \emph{An Introduction to Quantum Field Theory}, Addison-Wesley, Reading (1995).
\bibitem{Pomeranchuk45} I. Pomeranchuk and Ya. Smorodinsky, J. Phys. (USSR) \textbf{9}, 97 (1945).
\bibitem{Rau71} A. R. P. Rau, Phys. Rev. A \textbf{4}, 207 (1971).
\bibitem{Reinhardt77} J. Reinhardt and W. Greiner, Rep. Prog. Phys. \textbf{40}, 219 (1977).
\bibitem{Sadeghpour00} H. R. Sadeghpour, J. L. Bohn, M. J. Cavagnero, B. D. Esry, I. I. Fabrikant, J. H. Macek and A. R. P. Rau, J. Phys. B \textbf{33}, R93 (2000).
\bibitem{Sakharov48} A. D. Sakharov, Zh. Eksp. Teor. Fiz. \textbf{18}, 631 (1948).
\bibitem{Sauter31} F. Sauter, Z. Phys. \textbf{69}, 742 (1931).
\bibitem{Schwinger51} J. Schwinger, Phys. Rev. \textbf{82}, 664 (1951); \textbf{91}, 728 (1953).
\bibitem{Seaton58} M. J. Seaton, Mon. Not. R. Astron. Soc. \textbf{118}, 504 (1958).
\bibitem{Seaton83} M. J. Seaton, Rep. Prog. Phys. \textbf{46}, 167 (1983).
\bibitem{Sigmund} P. Sigmund, \emph{Particle Penetration and Radiation Effects}, Springer, Berlin (2006).
\bibitem{Sommerfeld31} A. Sommerfeld, Ann. Phys. (Leipzig) \textbf{403}, 257 (1931).
\bibitem{Stueckelberg32} E. C. G. Stueckelberg, Helv. Phys. Acta \textbf{5}, 369 (1932).
\bibitem{Taylor} J. R. Taylor, \emph{Scattering Theory: The Quantum Theory of Nonrelativistic Collisions}, Wiley, New York (1972); Dover reprint (2006).
\bibitem{Thomson12} J. J. Thomson, Phil. Mag. \textbf{23}, 449 (1912).
\bibitem{Wannier53} G. H. Wannier, Phys. Rev. \textbf{90}, 817 (1953).
\bibitem{Weinberg} S. Weinberg, \emph{The Quantum Theory of Fields}, vol.~1, Cambridge University Press (1995). Chapter 3 treats scattering theory and Chapter 13 infrared effects.
\bibitem{Wigner48} E. P. Wigner, Phys. Rev. \textbf{73}, 1002 (1948).
\bibitem{ZP71} Ya. B. Zel'dovich and V. S. Popov, Usp. Fiz. Nauk \textbf{105}, 403 (1971) [Sov. Phys. Usp. \textbf{14}, 673 (1972)].
\bibitem{Zener32} C. Zener, Proc. R. Soc. A \textbf{137}, 696 (1932).
\bibitem{PDG} F. Takahashi \emph{et al.} (Particle Data Group), \emph{Review of Particle Physics}, Int. J. Mod. Phys. A \textbf{41}, 2630011 (2026), \texttt{https://pdg.lbl.gov}. The reviews on kinematics and on the passage of particles through matter are the relevant ones.
\end{thebibliography}

\end{document}


